\renewcommand{\chaptertitle}{The Roots of Computing in the Ancient World}

%  Make first page footer:
\fancypagestyle{firststyle}{%
\fancyhf{}% Clear header/footer
\fancyhead{}
\fancyfoot[L]{{\footnotesize
              %% We toggle between these:
              % Manuscript submitted for review.\\
              \authorname{}~\authorpatp{}.  ``\chaptertitle{}''.  {\it The Nock Book} (2025):  1–$x$.
              }}
}
%  Arrange subsequent pages:
\fancyhf{}
\fancyhead[LE]{{\urbitfont The Nock Book}}
\fancyhead[RO]{\chaptertitle{}}
\fancyfoot[LE,RO]{\thepage}

\chapter{\chaptertitle}

\thispagestyle{firststyle}

\begin{abstract}
Cuneiform and early computational concepts are explored in this chapter.
\end{abstract}

\begin{quote}
Machines are pure products of human intellect, constructed for man’s own purposes. They are nowhere to be found in the world of nature (as products of nature); yet the workings of the laws of nature find their purest expression in machines, purer than in any of the products of nature itself. The laws of nature work directly in machines, with an immediacy not to be found in the products of nature.

(Keiji Nishitani, {\it Religion and Nothingness}, 1982, pp.~129--130)
\end{quote}

There are many worthy raconteurs of the history of computing, such as \citet{Hofstadter1980} and \citet{Penrose2004}.  What distinguishes my approach to the history of computing from other tellings?  Computing has deep connexions to alchemy and true language, which I hope to demonstrate without assuming any Hermetic background on the reader's part.  Much like technology, computing is not merely a tool but a stance towards the world.  However, it does not entail technologism's supercession; computing is always both a lens and a kaleidoscope, through which many layers of reality can be seen simultaneously.  It is a way of seeing the world as a system of pure statements and acting to make those statements true.  It is both ontology and epistemology.

First, therefore, I must establish that computation is the apex of a long human search for a language of truth.  The origin of symbol itself in the σύμβολον, a token broken in half to infallibly mark identity, is one font.  Later ``symbol'' was the profession of faith, both an attempt to formulate truth and a sign of shared claims.  The primary source, however, is language itself, Hofstadter's ``strange loop'' of self-reference.  The first possible Nock formula is a crash, the void.  The second is self-reference, awakening.  Peirce would be proud of the progression.  We have tricked ourselves into thinking that computing is a late development, a fruitful elaboration of the end of mathematics.  In fact, it is the philosopher's stone, and it underlies every process in the world.  It has turned lead into gold.

The historical account which follows aims throughout at the goal of understanding what Nock is and why it matters as a protocol.  The reader will understand what mankind has been looking for, when we have caught various pieces of the puzzle, how these things came together in the mid-20th century, and how they have been refined into a pure machine.

Outline:
- Cuneiform, tabulation, calculation, and algorithm
- Solid-state indices:  abaci, tokens, etc.  Mechanism as embodied equation.
- Pure language and true statements (Llull, Leibniz, Wilkins)
- Generalization (Jacquard, Babbage)
- The universal machine (Turing, G\"{o}del, Zuse, von Neumann)
- Pure statements as Pages from the Book (Curry, Church)

\section{Tabulation}

\section{Solid-state indices}

The problem is that any calculation lives in someone's head or someone's claim.  This is a summary, but not the proof or solution.
